
\section{\trCriticalPoint、\trCriticalValue{}与\trSardTheorem}
\label{sec:critical-points-sard}

本节可参见例如 [39]。

\subsection{\trCriticalPoint}
\label{subsec:critical-points}

\trManifold{} $M$ 的一个点 $x$ 称为映射
\[
f:M\longrightarrow N
\]
（取值于\trManifold{} $N$）的一个\textbf{\trCriticalPoint}（critical point），如果线性映射 $T_{x}f$ 不具有极大的\trRank：
\begin{itemize}
\item 或者 $\dim M\le\dim N$ 而 $f$ 在 $x$ 处不是一个\trImmersion；
\item 或者 $\dim M\ge\dim N$ 而 $f$ 在 $x$ 处不是一个\trSubmersion。
\end{itemize}

设 $m$、$n$ 分别为 $M$、$N$ 的维数——记号按字母表次序各就各位。当 $m\le n$ 时，条件为 $\dim\operatorname{Ker}T_{x}f>0$；当 $m\ge n$ 时，则为 $\operatorname{rk}T_{x}f<n$。我们可以把这概括为：$x$ 是\trCriticalPoint，如果 $f$ 在 $x$ 处的余秩为正——这里采用余秩（corank）的特设定义，即 $\inf(m,n)-\operatorname{rk}T_{x}f$。
\noindent\textbf{例（基本例）。}当 $N=\RR$ 时，我们有 $n=1$，而点 $x$ 是\trCriticalPoint{}当且仅当 $T_{x}f\equiv 0$。

\subsection{\trSardTheorem}
\label{subsec:sard-theorem}

映射 $f$ 的一个\textbf{\trCriticalValue}（critical value）是目标\trManifold{} $N$ 的这样的点：它是某个\trCriticalPoint{}的像。$N$ 的其余各点都是\textbf{\trRegularValue}（regular value）。我们特别要指出：一个\trRegularValue{}未必是一个值，也就是说，未必含于 $f$ 的像之中。此外，还可能有非\trCriticalPoint，它们的像是\trCriticalValue。
\trSardTheorem{}说的是：在某些条件下，一个可微映射的几乎所有值都是\trRegular{}的（注意，“零测集”的概念、因而“几乎处处”的概念，在\trManifold{}上是\trWellDefined{}的）。更精确地说，我们有下述结果。

\begin{theorem}[\trSardTheorem]\label{thm:sard-theorem}\BookIndexDefinition{trSardTheorem}
设 $f:M\to N$ 是一个 $\cC^{\infty}$ 映射；那么 $f$ 的诸\trCriticalValue{}所成之集在 $N$ 中是零测的。
\end{theorem}

\begin{remark}\label{thm:m-less-n-measure-zero}
当 $M$ 的维数 $m$ 小于 $N$ 的维数 $n$ 时，$f(M)$ 本身在 $N$ 中就是零测的。这是“子空间 $\RR^{m}$ 在 $\RR^{n}$ 中零测”这一事实的一个并非全然平凡的推论。例如，在这种情形下，从 $M$ 到 $N$ 的一个可微映射永远不是\trSurjective{}的。因此，这条定理主要是在 $m\ge n$ 时才令人感兴趣（而且也更难证明）。
\end{remark}

关于证明可参见例如 [39,~p.~34ff.]。

\subsection{函数的\trRegular\trLevelSet、\trSublevelSet}
\label{subsec:regular-level-sets-sublevel-sets}

如果 $f:V\to\RR$ 是一个函数，且 $\alpha$ 是 $f$ 的一个\trRegularValue，那么由\trRegularValue{}定理（即\cref{thm:submanifolds-characterization} 的刻画~(2)），\trLevelSet{} $f^{-1}(\alpha)$ 是 $V$ 的一个余维数为 $1$ 的\trSubmanifold。而\textbf{\trSublevelSet}（sublevel set）$V^{\alpha}=f^{-1}\bigl(\,]-\infty,\alpha]\bigr)$ 则反过来是一个以 $f^{-1}(\alpha)$ 为边界的\trManifoldWithBoundary。
设 $f:(V,\partial V)\to W$ 是一个从\trManifoldWithBoundary{}到一个\trManifold{} $W$ 的 $\cC^{\infty}$ 映射。我们设 $W$ 的点 $w$ 对 $f$ 与 $f|_{\partial V}$ 二者而言都是\trRegularValue；那么 $f^{-1}(w)$ 是 $V$ 的一个\trSubmanifold，其边界为
\[
\partial f^{-1}(w)=\bigl(f|_{\partial V}\bigr)^{-1}(w)\text{。}
\]
